\scanpage{18}
% Pencil foliation: 16, upper right; 17, lower left.
Considérons $\Gamma=\underline{\operatorname{Aut}}_{\mathbb{Q}}(V_0)$,
groupe algébrique
sur $\mathbb{Q}$, et les deux hom
\diagram{\begin{tikzcd}
i_1,i_2:\mathbb{G}_{\mathrm m,\mathbb{C}}\arrow[r]&\Gamma_{\mathbb{C}}
\end{tikzcd}}
définissant la bigraduation de $V$
\deleted{\unsure{$\varphi(t)$}}
\[
 i(\lambda,\mu)\longrightarrow\lambda^p\mu^q\quad\text{sur }V^{pq}
\]
Soit $G$ le plus petit ss-groupe algébrique de \inserted{$\Gamma$}
tel que $G_{\mathbb{C}}$ contienne $\operatorname{Im}i_1,\operatorname{Im}i_2$ ;
\deleted{Alors} c'est aussi le ss-groupe algébrique de
$\Gamma_{\mathbb{C}}$ engendré par les $i_i^\sigma$, $\sigma$ parcourant
les \deleted{\unsure{les}} autom de $\mathbb{C}$.
C'est un groupe
algébrique connexe, appelé \uline{groupe de Mumford-Tate}
de la structure de Hodge.

\deleted{Pour toute
\unsure{représentation} de $G$ \deleted{par} \unsure{dans}
\unsure{un vectoriel} \deleted{\unsure{par hom}}
$M$ de rang fini sur $\mathbb{Q}$,}

Si $\Sigma$ est
une opération tensorielle en une variable,
de sorte que $G$ opère sur $\Sigma(\unsure{W_0|V_0})$,
on voit que les éléments de $\Sigma(\unsure{W_0|V_0})$ invariants
sous $G$ sont exactement ceux qui sont
invariants par \inserted{les $i_i^\sigma$} $i_1,i_2$, i.e. qui sont de
bidegré $0,0$. Cela permet de voir que
